\documentclass[11pt]{article}

% ---------------------------------------------------------------
%  OCIA Student Handout
%  Order of Christian Initiation of Adults
% ---------------------------------------------------------------

\usepackage[utf8]{inputenc}
\usepackage[T1]{fontenc}
\usepackage[margin=1in]{geometry}
\usepackage{parskip}
\usepackage{fancyhdr}
\usepackage{xcolor}
\usepackage[expansion=false]{microtype}
\usepackage{multicol}
\usepackage[hidelinks]{hyperref}

% ---- Colours ----
\definecolor{churchblue}{RGB}{31,56,100}
\definecolor{churchgold}{RGB}{153,101,21}

% ---- Section styling (base LaTeX, no extra packages) ----
% Redefine the standard sectioning commands' fonts to add colour,
% using \@startsection so numbering and spacing remain standard.
\makeatletter
\renewcommand\section{\@startsection{section}{1}{\z@}%
  {-3.5ex \@plus -1ex \@minus -.2ex}%
  {2.3ex \@plus.2ex}%
  {\normalfont\Large\bfseries\color{churchblue}}}
\renewcommand\subsection{\@startsection{subsection}{2}{\z@}%
  {-3.25ex\@plus -1ex \@minus -.2ex}%
  {1.5ex \@plus .2ex}%
  {\normalfont\large\bfseries\color{churchgold}}}
\makeatother

% ---- Headers / footers ----
\pagestyle{fancy}
\fancyhf{}
\renewcommand{\headrulewidth}{0.4pt}
\lhead{\small\itshape OCIA Student Handout}
\rhead{\small\thepage}
\cfoot{}

% ---- Prayer environment (indented block) ----
\newenvironment{prayer}
  {\begin{quote}\itshape}
  {\end{quote}}

% ---------------------------------------------------------------
\begin{document}

% ---- Title page ----
\begin{titlepage}
  \centering
  \vspace*{2cm}
  {\Huge\bfseries\color{churchblue} A Handbook for OCIA Students\par}
  \vspace{0.8cm}
  {\Large Order of Christian Initiation of Adults\par}
  \vspace{1.5cm}
  {\large\itshape Prayers, Practices, and Teachings\\of the Catholic Faith\par}
  \vspace{2cm}
  {\large\color{churchgold}\rule{0.5\textwidth}{0.8pt}\par}
  \vspace{2cm}
  {\large \emph{``Go therefore and make disciples of all nations,}\\
   \emph{baptizing them in the name of the Father,}\\
   \emph{and of the Son, and of the Holy Spirit.''}\par}
  \vspace{0.4cm}
  {\normalsize --- Matthew 28:19\par}
  \vfill
  {\large For the instruction of catechumens and candidates\par}
  \vspace{0.5cm}
\end{titlepage}

\tableofcontents
\newpage

% ===============================================================
\section{Prayers}
% ===============================================================

Every Catholic should commit these foundational prayers to memory. They
form the daily vocabulary of our conversation with God.

% ---------------------------------------------------------------
\subsection*{i.~~The Our Father (The Lord's Prayer)}
\addcontentsline{toc}{subsection}{i. The Our Father}
\begin{prayer}
Our Father, who art in heaven,\\
hallowed be thy name;\\
thy kingdom come,\\
thy will be done\\
on earth as it is in heaven.\\
Give us this day our daily bread,\\
and forgive us our trespasses,\\
as we forgive those who trespass against us;\\
and lead us not into temptation,\\
but deliver us from evil.\\
Amen.
\end{prayer}

% ---------------------------------------------------------------
\subsection*{ii.~~The Hail Mary}
\addcontentsline{toc}{subsection}{ii. The Hail Mary}
\begin{prayer}
Hail Mary, full of grace,\\
the Lord is with thee.\\
Blessed art thou amongst women,\\
and blessed is the fruit of thy womb, Jesus.\\
Holy Mary, Mother of God,\\
pray for us sinners,\\
now and at the hour of our death.\\
Amen.
\end{prayer}

% ---------------------------------------------------------------
\subsection*{iii.~~The Glory Be (Doxology)}
\addcontentsline{toc}{subsection}{iii. The Glory Be}
\begin{prayer}
Glory be to the Father,\\
and to the Son,\\
and to the Holy Spirit.\\
As it was in the beginning,\\
is now, and ever shall be,\\
world without end.\\
Amen.
\end{prayer}

% ---------------------------------------------------------------
\subsection*{iv.~~The St. Michael Prayer}
\addcontentsline{toc}{subsection}{iv. The St. Michael Prayer}
\begin{prayer}
St. Michael the Archangel,\\
defend us in battle.\\
Be our protection against the wickedness and snares of the devil.\\
May God rebuke him, we humbly pray;\\
and do thou, O Prince of the heavenly host,\\
by the power of God,\\
cast into hell Satan and all the evil spirits\\
who prowl about the world seeking the ruin of souls.\\
Amen.
\end{prayer}

% ---------------------------------------------------------------
\subsection*{v.~~The Apostles' Creed}
\addcontentsline{toc}{subsection}{v. The Apostles' Creed}
\begin{prayer}
I believe in God,\\
the Father almighty,\\
Creator of heaven and earth,\\
and in Jesus Christ, his only Son, our Lord,\\
who was conceived by the Holy Spirit,\\
born of the Virgin Mary,\\
suffered under Pontius Pilate,\\
was crucified, died and was buried;\\
he descended into hell;\\
on the third day he rose again from the dead;\\
he ascended into heaven,\\
and is seated at the right hand of God the Father almighty;\\
from there he will come to judge the living and the dead.\\
I believe in the Holy Spirit,\\
the holy catholic Church,\\
the communion of saints,\\
the forgiveness of sins,\\
the resurrection of the body,\\
and life everlasting.\\
Amen.
\end{prayer}

% ---------------------------------------------------------------
\subsection*{vi.~~The Act of Contrition}
\addcontentsline{toc}{subsection}{vi. The Act of Contrition}
\begin{prayer}
O my God,\\
I am heartily sorry for having offended Thee,\\
and I detest all my sins\\
because of Thy just punishments,\\
but most of all because they offend Thee, my God,\\
who art all good and deserving of all my love.\\
I firmly resolve, with the help of Thy grace,\\
to sin no more\\
and to avoid the near occasions of sin.\\
Amen.
\end{prayer}

% ---------------------------------------------------------------
\subsection*{vii.~~The Anima Christi}
\addcontentsline{toc}{subsection}{vii. The Anima Christi}
\begin{prayer}
Soul of Christ, sanctify me.\\
Body of Christ, save me.\\
Blood of Christ, inebriate me.\\
Water from the side of Christ, wash me.\\
Passion of Christ, strengthen me.\\
O good Jesus, hear me.\\
Within Thy wounds hide me.\\
Suffer me not to be separated from Thee.\\
From the malignant enemy defend me.\\
In the hour of my death call me,\\
and bid me come unto Thee,\\
that with Thy saints I may praise Thee\\
forever and ever.\\
Amen.
\end{prayer}

\newpage
% ===============================================================
\section{Dress Code for Catholics in Church}
% ===============================================================

The way we dress for Mass and the sacraments is an outward sign of the
reverence we owe to God present in the Most Blessed Sacrament. Modesty,
dignity, and propriety should always guide our choices. ``Sunday best''
is the standard: we dress as though meeting the King of Kings, because
we are.

% ---------------------------------------------------------------
\subsection*{i.~~Attire for a Traditional (Tridentine) Mass}
\addcontentsline{toc}{subsection}{i. Traditional Mass Attire}

The same standard of reverence applies at every Mass, but the following
reflects the customary practice associated with the Traditional Latin
(Tridentine) Mass:

\textbf{For men:}
\begin{itemize}
  \item A suit, or at minimum dress trousers with a collared shirt; a jacket and tie are encouraged.
  \item Closed dress shoes; no athletic wear or sneakers.
  \item Hair neat and well-groomed.
\end{itemize}

\textbf{For women:}
\begin{itemize}
  \item A dress or skirt of modest length (covering the knee), or modest formal attire.
  \item Sleeves that cover the shoulders; necklines that are modest; backs covered.
  \item A chapel veil (mantilla) or other head covering is the traditional custom (see iii below).
  \item No low-cut, tight, or sheer clothing.
\end{itemize}

\textbf{General principles for all:}
\begin{itemize}
  \item Avoid shorts, jeans, T-shirts with slogans, beachwear, and casual sportswear.
  \item Clothing should be clean, neat, and not call undue attention to the wearer.
  \item Modesty covers both how much is shown and how the garment fits.
\end{itemize}

% ---------------------------------------------------------------
\subsection*{ii.~~What to Wear for Your Baptism}
\addcontentsline{toc}{subsection}{ii. Baptismal Attire}

Baptism clothes the newly baptized in Christ, symbolized by a
\textbf{white garment}. For adults received at the Easter Vigil or
otherwise:

\begin{itemize}
  \item \textbf{White is the colour of baptism}, signifying purity and the new life of grace. Where the parish provides a white robe or alb, wear it over your clothing.
  \item Beneath the garment, wear modest, dignified clothing as for Mass --- a white or light dress for women, a shirt and trousers for men.
  \item Choose clothing you do not mind getting wet, as water is poured (or you may be immersed) during the rite.
  \item Avoid heavy makeup or elaborate hairstyles that interfere with the pouring of water on the head.
  \item You will receive a white garment and a baptismal candle as part of the rite; treat both with reverence.
\end{itemize}

% ---------------------------------------------------------------
\subsection*{iii.~~Rules for Headwear (Men and Women)}
\addcontentsline{toc}{subsection}{iii. Headwear Rules}

The traditional discipline regarding head coverings is drawn from
St. Paul (1 Corinthians 11:4--7):

\textbf{For men:}
\begin{itemize}
  \item Men \textbf{remove} their hats and head coverings upon entering the church, as a sign of respect.
  \item Head coverings remain off throughout Mass and in the presence of the Blessed Sacrament.
  \item (Clergy wear liturgical head coverings, such as the biretta, as prescribed by the rubrics; this is distinct from lay practice.)
\end{itemize}

\textbf{For women:}
\begin{itemize}
  \item By long-standing custom, women \textbf{cover} their heads in church, traditionally with a chapel veil (mantilla) or hat.
  \item This practice, universal before 1983, remains a venerable and encouraged custom, especially at the Traditional Latin Mass.
  \item The veil is a sign of reverence and humility before God; it is worn from entering the church until leaving.
\end{itemize}

\textbf{Note:} Since the 1983 Code of Canon Law, head covering for women is
no longer a strict obligation but is widely retained as a pious and
beautiful tradition. Catechumens should follow the custom of their own
parish.

\newpage
% ===============================================================
\section{The Ten Commandments}
% ===============================================================

Given by God to Moses on Mount Sinai (Exodus 20:1--17; Deuteronomy
5:6--21), the Ten Commandments are the foundation of the moral law. The
Catholic enumeration is as follows:

\begin{enumerate}
  \item I am the \textsc{Lord} your God: you shall not have strange gods before me.
  \item You shall not take the name of the \textsc{Lord} your God in vain.
  \item Remember to keep holy the \textsc{Lord}'s Day.
  \item Honour your father and your mother.
  \item You shall not kill.
  \item You shall not commit adultery.
  \item You shall not steal.
  \item You shall not bear false witness against your neighbour.
  \item You shall not covet your neighbour's wife.
  \item You shall not covet your neighbour's goods.
\end{enumerate}

The first three commandments concern our duties toward God; the
remaining seven concern our duties toward our neighbour. Together they
are summarized by Our Lord in the two great commandments: to love God
above all things, and to love our neighbour as ourselves
(Matthew 22:37--40).

% ===============================================================
\section{The Five Precepts (Laws) of the Church}
% ===============================================================

The Precepts of the Church are positive laws that mark the minimum
required of the faithful in prayer and moral effort, in the growth of
love of God and neighbour (CCC 2041--2043):

\begin{enumerate}
  \item \textbf{To attend Mass on Sundays and Holy Days of Obligation}, and to rest from servile labour.
  \item \textbf{To confess your sins at least once a year} (and always before receiving Holy Communion in a state of mortal sin).
  \item \textbf{To receive Holy Communion at least during the Easter season} (the ``Easter duty'').
  \item \textbf{To observe the days of fasting and abstinence} established by the Church (e.g., Ash Wednesday and Good Friday; abstinence on Fridays).
  \item \textbf{To provide for the material needs of the Church}, each according to his ability.
\end{enumerate}

% ===============================================================
\section{Holy Days of Obligation (United States)}
% ===============================================================

In addition to every Sunday, the faithful in the United States are bound
to attend Mass on six Holy Days of Obligation. On these days, as on
Sundays, Catholics are to participate in Mass and to refrain from
unnecessary servile labour.

\begin{enumerate}
  \item \textbf{Solemnity of Mary, the Holy Mother of God} --- January~1
  \item \textbf{Ascension of the Lord} --- the Thursday after the Sixth Sunday of Easter (transferred to the Seventh Sunday of Easter in most U.S. dioceses)
  \item \textbf{Assumption of the Blessed Virgin Mary} --- August~15
  \item \textbf{All Saints} --- November~1
  \item \textbf{Immaculate Conception} of the Blessed Virgin Mary --- December~8 (patronal feast of the United States)
  \item \textbf{Nativity of Our Lord Jesus Christ} (Christmas) --- December~25
\end{enumerate}

\textbf{Note:} When the Solemnity of Mary Mother of God (Jan.~1), the
Assumption (Aug.~15), or All Saints (Nov.~1) falls on a Saturday or a
Monday, the obligation to attend Mass is abrogated in the United States.
The Immaculate Conception and Christmas are always days of obligation.

\subsection*{Dates in the Coming Year (September 2026 -- August 2027)}

The Holy Days of Obligation occurring in the twelve months beginning
September~1, 2026 fall on the following dates:

\begin{center}
\begin{tabular}{lll}
\textbf{Holy Day} & \textbf{Date} & \textbf{Day} \\[2pt]
\hline\\[-8pt]
All Saints & November 1, 2026 & Sunday \\
Immaculate Conception & December 8, 2026 & Tuesday \\
Nativity of Our Lord (Christmas) & December 25, 2026 & Friday \\
Mary, the Holy Mother of God & January 1, 2027 & Friday \\
Ascension of the Lord (Thursday) & May 6, 2027 & Thursday \\
\quad or, where transferred (Sunday) & May 9, 2027 & Sunday \\
Assumption of the B.V.M. & August 15, 2027 & Sunday \\
\end{tabular}
\end{center}

\textbf{Please note:} All Saints (Nov.~1, 2026) and the Assumption
(Aug.~15, 2027) fall on a Sunday and are celebrated at the Sunday Masses.
The Ascension is observed on Thursday, May~6, 2027 in the provinces that
retain Ascension Thursday, and on Sunday, May~9, 2027 in the dioceses
where it is transferred; consult your own parish.

\newpage
% ===============================================================
\section{The Seven Sacraments}
% ===============================================================

The sacraments are efficacious signs of grace, instituted by Christ and
entrusted to the Church, by which divine life is dispensed to us
(CCC 1131). They are grouped in three categories.

\subsection*{Sacraments of Initiation}
\begin{enumerate}
  \item \textbf{Baptism} --- Cleanses original sin and all personal sin, makes us children of God and members of the Church, and imprints an indelible character on the soul.
  \item \textbf{Confirmation} --- Completes baptismal grace; strengthens us with the gifts of the Holy Spirit to profess and defend the faith.
  \item \textbf{The Holy Eucharist} --- The true Body, Blood, Soul, and Divinity of Christ, the source and summit of the Christian life.
\end{enumerate}

\subsection*{Sacraments of Healing}
\begin{enumerate}\setcounter{enumi}{3}
  \item \textbf{Penance (Confession / Reconciliation)} --- Forgives sins committed after Baptism and reconciles the sinner with God and the Church.
  \item \textbf{Anointing of the Sick} --- Confers grace, comfort, and (when God wills) healing upon those gravely ill or near death.
\end{enumerate}

\subsection*{Sacraments at the Service of Communion}
\begin{enumerate}\setcounter{enumi}{5}
  \item \textbf{Holy Orders} --- Consecrates men as bishops, priests, and deacons to serve and govern the Church.
  \item \textbf{Matrimony} --- Unites a man and a woman in a lifelong, faithful, fruitful covenant, signifying the union of Christ and his Church.
\end{enumerate}

% ===============================================================
\section{Rules for Catholic Marriage}
% ===============================================================

Marriage in the Catholic Church is a sacred covenant and a sacrament. For
a valid and licit marriage, the Church requires that certain conditions
be met and that no impediments stand in the way.

\subsection*{General Requirements}
\begin{itemize}
  \item \textbf{Freedom to marry:} Both parties must be free of any prior valid marriage bond. (A previous marriage requires a declaration of nullity or proof of the death of the former spouse.)
  \item \textbf{Free consent:} Marriage is created by the free, mutual consent of the spouses, without coercion, fear, or fraud.
  \item \textbf{The three goods (essential properties):} The couple must intend \emph{permanence} (until death), \emph{fidelity} (exclusivity), and \emph{openness to children} (the procreation and education of offspring). Deliberately excluding any of these renders consent invalid.
  \item \textbf{Canonical form:} The marriage must normally be celebrated before a bishop, priest, or deacon and two witnesses (unless a dispensation is granted).
  \item \textbf{Proper preparation:} Couples must complete marriage preparation (pre-Cana) as required by the diocese.
  \item \textbf{Minimum age and capacity:} The parties must have attained the canonical age (16 for men, 14 for women in universal law, though civil and diocesan norms are typically higher) and possess sufficient use of reason.
\end{itemize}

\subsection*{Impediments to Marriage}
The following are among the diriment impediments that render a marriage
invalid (Code of Canon Law, cc. 1083--1094):
\begin{itemize}
  \item A prior valid marriage bond.
  \item Impotence (antecedent and perpetual).
  \item Disparity of cult (one party unbaptized) without dispensation.
  \item Holy Orders, or a public perpetual vow of chastity.
  \item Abduction, crime (conjugicide), or coercion.
  \item Relationships of consanguinity, affinity, or legal/adoptive kinship (see below).
\end{itemize}

\subsection*{Consanguinity Rules (Blood Relatives)}

Consanguinity is the relationship between persons descended from a common
ancestor. The Church forbids marriage between close blood relatives, both
for moral and biological reasons (CIC c. 1091):

\begin{itemize}
  \item \textbf{Direct line:} Marriage is \emph{always invalid} between all
        ancestors and descendants --- whether legitimate or natural. (For
        example, a parent and child, or grandparent and grandchild, can
        \emph{never} marry, and no dispensation is possible.)
  \item \textbf{Collateral line:} Marriage is invalid up to and including
        the \emph{fourth degree} of the collateral line. This means:
        \begin{itemize}
          \item \textbf{2nd degree} --- brother and sister: \emph{never} permitted (no dispensation).
          \item \textbf{3rd degree} --- uncle/aunt with niece/nephew: forbidden; rarely dispensed.
          \item \textbf{4th degree} --- first cousins: forbidden, but a dispensation may be granted by the bishop for a just cause.
        \end{itemize}
  \item \textbf{Beyond the fourth degree}, consanguinity is not an impediment (e.g., second cousins may marry without dispensation).
\end{itemize}

\textbf{Related impediments:}
\begin{itemize}
  \item \textbf{Affinity} (c. 1092): the relationship arising from a valid
        marriage between a spouse and the blood relatives of the other
        spouse. Affinity in the direct line (e.g., a man and his deceased
        wife's mother or daughter) invalidates marriage.
  \item \textbf{Adoptive / legal kinship} (c. 1094): those related by legal
        adoption in the direct line, or in the second degree of the
        collateral line (adoptive siblings), cannot validly marry.
\end{itemize}

\vspace{1em}
\begin{center}
\textit{When in doubt about any impediment or requirement, always consult
your parish priest, who will guide you and seek any necessary
dispensation from the diocese.}
\end{center}

\vfill
\begin{center}
\rule{0.4\textwidth}{0.6pt}\\
\vspace{0.4em}
\small This handout is provided for catechetical instruction.\\
For authoritative reference, consult the \emph{Catechism of the Catholic Church}\\
and the \emph{Code of Canon Law}.
\end{center}

\end{document}
